\section{پاسخ سوال ۷}
\begin{stepbox}
\field{تفسیر طیف تبدیل فوریه \lr{(Fourier Transform)}}{
تصویر اصلی سمت چپ، یک الگوی بافت‌دار (Texture) است که از مربع‌های روشن و تاریک تشکیل شده و به صورت یک \textbf{شبکه تکرارشونده با زاویه ۴۵ درجه} (شبیه شطرنجی مورب) قرار گرفته است.

مکان ضرایب بزرگ در طیف تبدیل فوریه به دلایل زیر توجیه می‌شود:
\begin{enumerate}
    \item \textbf{جهت‌گیری (Orientation):} چون الگوی فضایی اصلی تصویر دارای لبه‌ها و تناوب در امتداد قطرهای ۴۵ درجه و ۱۳۵ درجه است، انرژی فرکانسی در دامنه فوریه نیز در امتداد خطوطی با همین زوایا (خطوط عمود بر لبه‌ها) متمرکز می‌شود. به همین دلیل شکل یک تقاطع ستاره‌ای (صلیبی مورب) در طیف دیده می‌شود.
    \item \textbf{فرکانس پایه و هارمونیک‌ها (Frequency \& Harmonics):} لکه‌های روشن (ضرایب بزرگ) روی این محورهای قطری نشان‌دهنده فرکانس‌های تکرار الگوها هستند. لکه‌های نزدیک‌تر به مرکز مختصات (مبدا DC)، معادل فرکانس پایه (فاصله بین مربع‌های بزرگ) هستند. لکه‌های دورتر، \textbf{هارمونیک‌های فرد} فرکانس پایه می‌باشند. به دلیل تیز بودن لبه‌های مربع‌ها (تبدیل شکل موج مربعی به سینوسی‌ها)، طیف فرکانسی نیازمند هارمونیک‌های بالا برای ساختن لبه‌های تیز در حوزه مکان است.
\end{enumerate}
}
\end{stepbox}
